\chapter{Đạo Sống Qua Muôn Đời (The Way of Life Across the Ages)}
\begin{center}
  \textit{\color{truyengold}``Trăm năm trồng người --- con đường ấy chưa bao giờ ngắn.''}\\[4pt]
  \textit{(A hundred years to raise a person --- that road has never been short.)}\\[4pt]
  \small\color{truyenblue}--- Quản Trọng ---
\end{center}
\ngancach

\section{Quan Vũ Thả Tào Tháo Ở Hoa Dung Đạo}
\begin{truyen}{Quan Vũ Thả Tào Tháo Ở Hoa Dung Đạo}{Tam Quốc Chí --- Trung Quốc, thế kỷ 3}
\chuhoa{S}{au} trận Xích Bích, Tào Tháo bại trận tháo chạy. Quan Vũ được lệnh chặn đường ở Hoa Dung Đạo; nếu bắt được Tào Tháo, chiến công sẽ là vô song.\\[4pt]
\textit{(After the Battle of Red Cliffs, Cao Cao fled in defeat. Guan Yu was ordered to block the road at Huarong Path; capturing Cao Cao would be an unmatched achievement.)}

Tào Tháo quỳ xuống nhắc lại ân xưa: năm xưa Quan Vũ bị vây hãm, chính Tào Tháo đã cứu ông thoát.\\[4pt]
\textit{(Cao Cao knelt and recalled past kindness: years before, when Guan Yu was surrounded, it was Cao Cao who saved him.)}

Quan Vũ biết nếu thả Tào Tháo ông sẽ bị tội chết theo quân luật --- nhưng ông không thể để người từng có ơn với mình bị giết chết dưới tay mình.\\[4pt]
\textit{(Guan Yu knew that releasing Cao Cao meant facing execution under military law --- but he could not let the man who once showed him kindness die by his hand.)}

Ông bước sang một bên, để Tào Tháo và tàn quân đi qua.\\[4pt]
\textit{(He stepped aside, allowing Cao Cao and his remnant troops to pass.)}

\end{truyen}
\begin{baihoc}
  Nghĩa khí đôi khi đòi hỏi ta phải chọn lương tâm thay vì lợi ích --- đó là dấu hiệu của người quân tử.\\[4pt]
  \textit{(Righteousness sometimes demands we choose conscience over benefit --- that is the mark of a person of integrity.)}
\end{baihoc}

\section{Thầy Chu Văn An Dâng Thất Trảm Sớ}
\begin{truyen}{Thầy Chu Văn An Dâng Thất Trảm Sớ}{Chu Văn An --- Việt Nam, thế kỷ 14}
\chuhoa{D}{ưới} thời Trần Dụ Tông, triều đình lộng hành, bảy kẻ gian thần hoành hành nước nhà. Chu Văn An --- bậc thầy của cả vua --- dâng sớ xin chém đầu bảy tên đó.\\[4pt]
\textit{(Under King Tran Du Tong, the court was corrupt, seven wicked ministers wreaking havoc on the realm. Chu Van An --- teacher even to the king --- submitted a memorial requesting the execution of those seven.)}

Vua không nghe. Chu Văn An lập tức từ quan, về Chí Linh dạy học và viết sách đến cuối đời.\\[4pt]
\textit{(The king refused. Chu Van An immediately resigned, returned to Chi Linh to teach and write until the end of his life.)}

Ông không hận, không thù, không cúi đầu --- chỉ âm thầm tiếp tục trồng người cho đất nước.\\[4pt]
\textit{(He harbored no resentment, no hatred, no bowed head --- quietly continuing to nurture the country's future people.)}

Hậu thế phong ông là ``Vạn thế sư biểu'' --- Người thầy của muôn đời.\\[4pt]
\textit{(Posterity honored him as `Teacher of Ten Thousand Generations'.)}

\end{truyen}
\begin{baihoc}
  Người thầy chân chính không chỉ dạy bài học --- mà sống như một bài học.\\[4pt]
  \textit{(A true teacher does not only teach lessons --- but lives as a lesson.)}
\end{baihoc}

\section{Nelson Mandela Tha Thứ Kẻ Cai Ngục}
\begin{truyen}{Nelson Mandela Tha Thứ Kẻ Cai Ngục}{Nelson Mandela --- Nam Phi, 1994}
\chuhoa{S}{au} 27 năm tù đày ở Robben Island, khi Nelson Mandela bước ra để tuyên thệ nhậm chức Tổng thống Nam Phi, ông mời ba người cai ngục từng giam cầm ông ngồi ở hàng ghế danh dự.\\[4pt]
\textit{(After 27 years imprisoned on Robben Island, when Nelson Mandela stepped out to be sworn in as President of South Africa, he invited three of the prison guards who had once held him captive to sit in the seats of honor.)}

Ông nói: ``Khi tôi bước ra khỏi cổng nhà tù, tôi biết rằng nếu không để lại sự cay đắng và thù hận phía sau, tôi vẫn còn là tù nhân.''\\[4pt]
\textit{(He said: `When I walked out the prison gate, I knew that if I did not leave my bitterness and hatred behind, I would still be a prisoner.')}

Mandela lựa chọn tha thứ không phải vì kẻ thù xứng đáng được tha thứ --- mà vì chính ông xứng đáng được tự do.\\[4pt]
\textit{(Mandela chose forgiveness not because his enemies deserved it --- but because he himself deserved freedom.)}

\end{truyen}
\begin{baihoc}
  Tha thứ là hành động dũng cảm nhất --- vì nó giải phóng người tha thứ, không phải người được tha thứ.\\[4pt]
  \textit{(Forgiveness is the most courageous act --- for it frees the one who forgives, not the one who is forgiven.)}
\end{baihoc}

\section{Lão Tử Và Triết Lý Nước}
\begin{truyen}{Lão Tử Và Triết Lý Nước}{Đạo Đức Kinh --- Trung Quốc, thế kỷ 6 TCN}
\chuhoa{L}{ão} Tử dạy: ``Thượng thiện nhược thủy'' --- Điều tốt đẹp nhất giống như nước.\\[4pt]
\textit{(Laozi taught: `The highest good is like water' --- Shang shan ruo shui.)}

Nước có lợi cho vạn vật mà không tranh giành; nó chảy xuống chỗ thấp mà người đời không ai muốn ở.\\[4pt]
\textit{(Water benefits all things without competing; it flows to low places that no one desires.)}

Nước mềm nhất thế gian nhưng chảy mãi thì khoét được đá cứng nhất.\\[4pt]
\textit{(Water is the softest thing in the world, yet flowing steadily it carves the hardest stone.)}

Người biết đạo sống như nước: khiêm nhường mà có sức mạnh vô song; nhường chỗ thấp mà cuối cùng lại chảy ra biển cả.\\[4pt]
\textit{(The one who understands the Way lives like water: humble yet possessing unmatched power; yielding to low places yet ultimately flowing to the sea.)}

\end{truyen}
\begin{baihoc}
  Sức mạnh thực sự không nằm ở sự cứng rắn --- mà ở sự kiên trì, mềm mại và không ngừng chảy về phía trước.\\[4pt]
  \textit{(True strength lies not in rigidity --- but in persistence, gentleness, and ceaseless forward flow.)}
\end{baihoc}
